(a) We have $D(f)=\varnothing$ exactly when $f$ lies in every prime ideal, that is, when $f$ is nilpotent.

(b) On a distinguished open $D(a)\subseteq X$, the map $f^\#$ is $A_a\to B_{\varphi(a)}$. If $\varphi$ is injective, these maps are injective by exactness of localization, so $f^\#$ is injective. Conversely, injectivity on global sections gives injectivity of $\varphi$. If $D(a)$ is disjoint from $f(Y)$, then $D(\varphi(a))=\varnothing$, so $\varphi(a)$ is nilpotent. Injectivity makes $a$ nilpotent, hence $D(a)=\varnothing$. Thus every nonempty open meets $f(Y)$.

(c) Identify $B=A/I$. Contraction identifies $\operatorname{Spec}B$ homeomorphically with $V(I)$. At a point $\mathfrak p\in V(I)$, the map on stalks of $\mathcal O_X\to f_*\mathcal O_Y$ is the surjection $A_{\mathfrak p}\to(A/I)_{\mathfrak p/I}$. Outside $V(I)$ its target stalk is $0$. Hence $f^\#$ is surjective.

(d) Factor $f$ as $Y\xrightarrow{h}X'=\operatorname{Spec}(A/\ker\varphi)\xrightarrow{i}X$. By (b), $h$ is dominant and $h^\#$ is injective; by (c), $i$ is a closed immersion. Since $f$ is a homeomorphism onto a closed subset, $h$ is a homeomorphism onto a closed subset of $X'$, and dominance makes it a homeomorphism onto $X'$. Under this identification, surjectivity of $f^\#=i_*h^\#\circ i^\#$ implies surjectivity of $h^\#$ on stalks. Thus $h$ is an isomorphism. Taking global sections gives $A/\ker\varphi\cong B$, so $\varphi$ is surjective.
